%%%%%%%%%%%%%%%%%%%%%%%%%%%%%%%%%%%%%%%%%
% Cleese Assignment (For Students)
% LaTeX Template
% Version 2.0 (27/5/2018)
%
% This template originates from:
% http://www.LaTeXTemplates.com
%
% Author:
% Vel (vel@LaTeXTemplates.com)
%
% License:
% CC BY-NC-SA 3.0 (http://creativecommons.org/licenses/by-nc-sa/3.0/)
% 
%%%%%%%%%%%%%%%%%%%%%%%%%%%%%%%%%%%%%%%%%

%----------------------------------------------------------------------------------------
%	PACKAGES AND OTHER DOCUMENT CONFIGURATIONS
%----------------------------------------------------------------------------------------

\documentclass[10pt]{article}
\usepackage{hyperref}
%%%%%%%%%%%%%%%%%%%%%%%%%%%%%%%%%%%%%%%%%
% Cleese Assignment
% Structure Specification File
% Version 1.0 (27/5/2018)
%
% This template originates from:
% http://www.LaTeXTemplates.com
%
% Author:
% Vel (vel@LaTeXTemplates.com)
%
% License:
% CC BY-NC-SA 3.0 (http://creativecommons.org/licenses/by-nc-sa/3.0/)
% 
%%%%%%%%%%%%%%%%%%%%%%%%%%%%%%%%%%%%%%%%%

%----------------------------------------------------------------------------------------
%	PACKAGES AND OTHER DOCUMENT CONFIGURATIONS
%----------------------------------------------------------------------------------------

\usepackage{lastpage} % Required to determine the last page number for the footer

\usepackage{graphicx} % Required to insert images

\setlength\parindent{0pt} % Removes all indentation from paragraphs

\usepackage[svgnames]{xcolor} % Enabling colors by their 'svgnames'
\usepackage[most]{tcolorbox} % Required for boxes that split across pages

\usepackage{booktabs} % Required for better horizontal rules in tables

\usepackage{listings} % Required for insertion of code

\usepackage{etoolbox} % Required for if statements

\usepackage{multido} % required for dotted lines

\usepackage[french]{babel} % English language hyphenation

\usepackage{svg}

%----------------------------------------------------------------------------------------
%	MARGINS
%----------------------------------------------------------------------------------------

\usepackage{geometry} % Required for adjusting page dimensions and margins

\geometry{
	paper=a4paper, % Change to letterpaper for US letter
	top=3cm, % Top margin
	bottom=3cm, % Bottom margin
	left=2.5cm, % Left margin
	right=2.5cm, % Right margin
	headheight=14pt, % Header height
	footskip=1.4cm, % Space from the bottom margin to the baseline of the footer
	headsep=1.2cm, % Space from the top margin to the baseline of the header
	%showframe, % Uncomment to show how the type block is set on the page
}

%----------------------------------------------------------------------------------------
%	FONT
%----------------------------------------------------------------------------------------
\usepackage{fontspec}
\usepackage{unicode-math}
\usepackage[utf8]{inputenc} % Required for inputting international characters
\usepackage[T1]{fontenc} % Output font encoding for international characters

\setmainfont[Ligatures=TeX]{Caladea}
%----------------------------------------------------------------------------------------
%	HEADERS AND FOOTERS
%----------------------------------------------------------------------------------------

\usepackage{fancyhdr} % Required for customising headers and footers

\pagestyle{fancy} % Enable custom headers and footers

% \lhead{\small\assignmentClass\ifdef{\assignmentClassInstructor}{\ (\assignmentClassInstructor):}{}\ \assignmentTitle} % Left header; output the instructor in brackets if one was set
% \chead{} % Centre header
% \rhead{\small\ifdef{\assignmentAuthorName}{\assignmentAuthorName}{\ifdef{\assignmentDueDate}{Due\ \assignmentDueDate}{}}} % Right header; output the author name if one was set, otherwise the due date if that was set

% \lfoot{} % Left footer
% \cfoot{\small Page\ \thepage\ sur\ \pageref{LastPage}} % Centre footer
% \rfoot{} % Right footer

% \renewcommand\headrulewidth{0.5pt} % Thickness of the header rule


% \fancypagestyle{plain}{
% 	\lhead{} % Left header; output the instructor in brackets if one was set
% 	\chead{} % Centre header
% 	\rhead{}

%   \renewcommand{\headrulewidth}{0pt}%
%   \fancyfoot[C]{\small Page\ \thepage\ sur\ \pageref{LastPage}}%
% }

%----------------------------------------------------------------------------------------
%	MODIFY SECTION STYLES
%----------------------------------------------------------------------------------------

% \usepackage{titlesec} % Required for modifying sections

%------------------------------------------------
% Section

% \titleformat
% {\section} % Section type being modified
% [block] % Shape type, can be: hang, block, display, runin, leftmargin, rightmargin, drop, wrap, frame
% {\Large\bfseries} % Format of the whole section
% {\assignmentQuestionName~\thesection} % Format of the section label
% {6pt} % Space between the title and label
% {} % Code before the label

% \titlespacing{\section}{0pt}{0.5\baselineskip}{0.5\baselineskip} % Spacing around section titles, the order is: left, before and after

%------------------------------------------------
% Subsection

% \titleformat
% {\subsection} % Section type being modified
% [block] % Shape type, can be: hang, block, display, runin, leftmargin, rightmargin, drop, wrap, frame
% {\itshape} % Format of the whole section
% {(\alph{subsection})} % Format of the section label
% {4pt} % Space between the title and label
% {} % Code before the label

% \titlespacing{\subsection}{0pt}{0.5\baselineskip}{0.5\baselineskip} % Spacing around section titles, the order is: left, before and after

% \renewcommand\thesubsection{(\alph{subsection})}

%----------------------------------------------------------------------------------------
%	CUSTOM QUESTION COMMANDS/ENVIRONMENTS
%----------------------------------------------------------------------------------------


%	TITLE SECTION
%----------------------------------------------------------------------------------------

% \newcommand{\authorstyle}[1]{{\large\usefont{OT1}{phv}{b}{n}#1}} % Authors style (Helvetica)

% \newcommand{\institution}[1]{{\footnotesize\usefont{OT1}{phv}{m}{sl}#1}} % Institutions style (Helvetica)

% \usepackage{titling} % Allows custom title configuration

% \newcommand{\HorRule}{\rule{\linewidth}{1pt}} % Defines the gold horizontal rule around the title

% \pretitle{
% 	\centering
% 	\vspace{-120pt} % Move the entire title section up
% 	\HorRule\vspace{10pt} % Horizontal rule before the title
% 	% \textbf
% 	\bfseries
% 	\fontsize{32}{36}
% 	\usefont{T1}{phv}{b}{n}
% 	\selectfont % Helvetica
% 	\color{DarkRed} % Text colour for the title and author(s)
% }

% \posttitle{\par\vskip 15pt} % Whitespace under the title

% \preauthor{} % Anything that will appear before \author is printed

% \postauthor{ % Anything that will appear after \author is printed
% 	% \vspace{10pt} % Space before the rule
% 	\par\HorRule % Horizontal rule after the title
% 	\vspace{-30pt} % Space after the title section
% }

 % Include the file specifying the document structure and custom commands

%----------------------------------------------------------------------------------------
%	ASSIGNMENT INFORMATION
%----------------------------------------------------------------------------------------

% Required
\newcommand{\assignmentQuestionName}{Question} % The word to be used as a prefix to question numbers; example alternatives: Problem, Exercise
\newcommand{\assignmentClass}{Physique Chimie} % Course/class
\newcommand{\assignmentTitle}{Activité n°1} % Assignment title or name
\newcommand{\assignmentAuthorName}{Chapitre 1} %

%----------------------------------------------------------------------------------------
%	VARIABLES
%----------------------------------------------------------------------------------------

\newcommand{\titreActivite}{Activité 2 : L'énergie verte} % titre de l'activité
\newcommand{\objectif}{ 	
	
	\begin{itemize}
		\item Se familiariser avec l'économie d'énergie.
		\item Connaître et comprendre les différentes énergies renouvelables.
	\end{itemize}
} % titre de l'activité
\newcommand{\contexte}{}
\newcommand{\resumeContexte}{Aujourd'hui, découvrons ensemble les différents types d'énergie qui existent.} % titre de l'activité




%----------------------------------------------------------------------------------------

\begin{document}
%----------------------------------------------------------------------------------------
%	TITLE PAGE
%----------------------------------------------------------------------------------------
\date{}
\title{\titreActivite}
\maketitle % Print the title page

%----------------------------------------------------------------------------------------
%	QUESTION 1
%----------------------------------------------------------------------------------------

\underline{\textbf{Objectif}} :  \vspace{2pt}
\objectif

% \vspace{8pt}

\begin{center}
	\includegraphics[width=0.15\columnwidth]{activité.jpg} % Example image
\end{center}

\textbf{\resumeContexte}
\vspace{-12pt}
%% ----------------------------------

\assignmentSection{Votre mission travail}


\textbf{\color{DarkBlue} {Document 1 : Vidéo}} \url{
	https://www.youtube.com/watch?v=kCYyrkBjdv4
}

%----------------------------------------------------------------------------------------
%	QUESTIONs	
%----------------------------------------------------------------------------------------
Répondre aux questions sur ton cahier ou au dos de l'activité.

\begin{question}
	\questiontext{
		Que fait-on chaque jour qui consomme de l'énergie?
	}
\end{question}

	courses journal, douche bain chaud radiateur

    \begin{question}
	\questiontext{
		Combien un français consomme t-il de pétrole par ans?
	}
\end{question}

4tonnes

\begin{question}
	\questiontext{
		Combien toute les humains consomment-ils de pétrole par ans?
	}
\end{question}

9.7milliard de tonne de pétrole par planète par ans

\begin{question}
	\questiontext{
		Quel rapport entre le pétrole et le dioxyde de carbone?
	}
\end{question}

\begin{question}
	\questiontext{
		Quels sont le ou les désavantages de l'énergie nucléaire?
	}
\end{question}

déchet nucléaire

\begin{question}
	\questiontext{
		Quels sont les 3 types de sources d'énergie que les humains utilisent?
	}
\end{question}

\begin{question}
	\questiontext{
		Quel part de l'énergie produite en France représente le nucléaire?
	}
\end{question}

40\% en france

\begin{question}
	\questiontext{
		Combien (en pourcentage) les pays industrialisés consomment-ils d'énergie? 
	}
\end{question}

64\% energie consomme par les pays indistrualisés 

\begin{question}
	\questiontext{
		De combien de degré la température sur Terre risque d'augmenter d'ici un siècle?
	}
\end{question}

(1 à 6)

\begin{question}
	\questiontext{
		À partir de quel vitesse de vent une éolienne démarre t-elle?
	}
\end{question}

vitesse du vent pour élolieene qui démarre, 10 à 15kmh

\begin{question}
	\questiontext{
		Qu'est-ce qu'un Watt?
	}
\end{question}

Watt puissance d'un appareil ou celle necessaire pour le faire fonctionner

\begin{question}
	\questiontext{
		Combien d'heures par an l'éolienne produit elle 1.3 mégawattheure?
		}
\end{question}

2000heure par an
	
\begin{question}
	\questiontext{
		Combien d'éoliennes faudrait-il pour remplacer un réacteur nulcéaire?
		}
\end{question}

4000

\begin{question}
	\questiontext{
		Quels panneaux peut transformer directement la lumière du soleil en électricité? 
		}
	\end{question}

photovoltaïque

\begin{question}
	\questiontext{
		Quels machines il ne faut pas utiliser en même temps
		pour ne pas avoir un problème de puissance ?
		}
	\end{question}

    machine à laver + aspirateur + four

    \begin{question}
	\questiontext{
		Qu'est-ce qu'un courant electrique ?
		}
	\end{question}

    déplacement d'électron

    \begin{question}
	\questiontext{
		Quelle est la vitesse maximale d'un taxi à hydrogène et
		pendant combien de kilomètre peut-il rouler?
		}
	\end{question}

    vmax=95km h  pdt 250/300km

\begin{question}
	\questiontext{
		Combien d'éoliennes faudrait-il pour remplacer un réacteur nulcéaire?
		}
\end{question}

4000

\begin{question}
	\questiontext{
		D'après toi qu'est-ce qu'une énergie renouvelable?
		}
\end{question}


\end{document}
